
Memorization in deep neural networks refers to the phenomenon 
where a model stores specific training examples rather than 
learning generalizable patterns from the underlying data 
distribution. Formally, given a training set 
$\mathcal{D} = \{(x_i, y_i)\}_{i=1}^n$, a model $f_\theta$ 
is said to memorize a training example $(x_i, y_i)$ if its 
prediction is driven by instance-specific features rather 
than the true underlying distribution $p(y|x)$.

\subsubsection{Generalization and Memorization}

The classical generalization framework defines the 
generalization error as:

\begin{equation}
\mathcal{E}_{gen} = \mathcal{E}_{test} - \mathcal{E}_{train}
\label{eq:generror}
\end{equation}

Where $\mathcal{E}_{train}$ is the empirical training error 
and $\mathcal{E}_{test}$ is the test error. A model that 
memorizes training examples achieves low $\mathcal{E}_{train}$ 
while exhibiting high $\mathcal{E}_{test}$, resulting in a 
large generalization gap. However, as demonstrated by 
\cite{chiyuan_zhang_bd8f0549}, this classical view is insufficient 
to fully explain memorization — deep neural networks can 
achieve near-zero training error even on randomly labeled 
data, yet still generalize well on correctly labeled data. 
This paradox motivates a more refined formalization.

\subsubsection{Memorization Score}

To quantify memorization at the instance level, the 
memorization score of a training example $(x_i, y_i)$ 
is defined as the influence of that example on the 
model's prediction \cite{feldman2021doeslearningrequirememorization}:

\begin{equation}
\text{mem}(f, \mathcal{D}, i) = 
\Pr_{f \sim \mathcal{A}(\mathcal{D})}
[f(x_i) = y_i] - 
\Pr_{f \sim \mathcal{A}(\mathcal{D} \setminus \{i\})}
[f(x_i) = y_i]
\label{eq:memscore}
\end{equation}

Where $\mathcal{A}$ is the training algorithm, 
$\mathcal{D} \setminus \{i\}$ denotes the training set 
with example $(x_i,y_i)$ removed, and the difference measures 
how much the model's correct prediction of $x_i$ depends 
on having seen that example during training. A high 
memorization score indicates that the model has memorized 
the example rather than learned a generalizable pattern.

\subsubsection{Label Memorization}

Label memorization occurs when the model fits the training 
labels without learning the underlying feature distributions. 
Formally, given a corrupted dataset 
$\tilde{\mathcal{D}} = \{(x_i, \tilde{y}_i)\}_{i=1}^n$ 
where $\tilde{y}_i$ is a randomly assigned label, a model 
exhibits label memorization if:

\begin{equation}
\mathcal{E}_{train}(\tilde{\mathcal{D}}) \approx 0 
\quad \text{while} \quad 
\mathcal{E}_{test}(\tilde{\mathcal{D}}) \approx 
\frac{K-1}{K}
\label{eq:labelmem}
\end{equation}

Where $K$ is the number of classes. This result, 
demonstrated empirically by \cite{chiyuan_zhang_bd8f0549} on 
CIFAR-10, shows that the model achieves near-zero 
training error on randomly labeled data while test 
error approaches the random chance baseline of 
$\frac{K-1}{K}$.

\subsubsection{Feature Memorization}

Feature memorization refers to the model retaining 
spurious features from specific training instances 
rather than abstracting meaningful representations 
\cite{simone_bombari_23e6de9b,yang_yu_4720d432}. Unlike 
label memorization, which concerns the retention of 
incorrect class assignments, feature memorization 
reflects the model's tendency to rely on instance-level 
surface features uncorrelated with the true underlying 
data distribution. Formally, let $x_i = v_c + v_s + \xi_i$ 
where $v_c$ is the core feature, $v_s$ is the spurious 
feature, and $\xi_i$ is instance-level noise. A model 
exhibits feature memorization when:

\begin{equation}
|f(v_s; W_T, z_T)| \gg |f(v_c; W_T, z_T)|
\label{eq:featuremem}
\end{equation}

That is, the model's output is predominantly determined 
by the spurious feature $v_s$ rather than the core 
feature $v_c$ \cite{yang_yu_4720d432}.